% !TeX root = ../../Book3.tex
% 纲要标题：### 6.3 强零点定理
% 纲要位置：工作记录/计划纲要.md，第 2076 行。
% 纲要细目（写作范围）：
% * 根理想与零点集的双向对应
% * Rabinowitsch 技巧；任意交换含幺环上的增广理想判据
% * 几何对象与代数对象反变关系的起点；普通坐标环的约化性

\section{强零点定理}\label{b3:sec:0603}

弱形式只判断方程组是否有解。强形式进一步回答：已知全部解之后，究竟能恢复哪些多项式关系？一个多项式与它的正整数次幂有相同零点，因此期望恢复原理想通常过强；正确对象是根理想。本节除特别说明外，\(k\) 始终为代数闭域。

\subsection{根理想与零点集的对应}\label{b3:subsec:0603:01}

设 \(k\) 为域，\(n\ge0\)，\(X\subseteq k^n\)。定义 \(X\) 的\term[lingdian-lixiang]{零点理想}[vanishing ideal]
\[
I(X)=\bigl\{\,f\in k[x_1,\ldots,x_n]\mid f(a)=0\text{ 对每个 }a\in X\,\bigr\}.
\]
这是根理想，因为 \(f(a)^r=0\) 在域中推出 \(f(a)=0\)。空集的零点理想按全称命题约定为整个多项式环。
\symidx{\(I(X)\)：点集 \(X\) 上恒为零的多项式组成的理想}

\begin{theorem}[Hilbert 强零点定理]\label{b3:thm:strong-nullstellensatz}
设 \(k\) 为代数闭域，\(n\ge0\)，\(J\) 为 \(k[x_1,\ldots,x_n]\) 的理想。记 \(Z_k(J)=\{a\in k^n:f(a)=0\text{ 对所有 }f\in J\}\)；对 \(X\subseteq k^n\)，记 \(I(X)=\{f\in k[x_1,\ldots,x_n]:f(a)=0\text{ 对所有 }a\in X\}\)。则
\[
I\bigl(Z_k(J)\bigr)=\sqrt J.
\]
这称为\term[qiang-lingdian-dingli]{强零点定理}[strong Nullstellensatz]。因此根理想与形如 \(Z_k(J)\) 的代数集，分别经 \(Z_k\) 与 \(I\) 构成反向包含的一一对应。
\end{theorem}

例如 \((x^2,y)\) 和 \((x,y)\) 在 \(k^2\) 中都只给出原点；从这个点集恢复的是根理想 \((x,y)\)。商环 \(k[x,y]/(x^2,y)\) 中保留的非零幂零元，无法由普通点集上的取值看到。下一小节证明等式，再核对对应关系。

\subsection{Rabinowitsch 技巧}\label{b3:subsec:0603:02}

把“不等于零”的条件写成“存在倒数”，是从弱形式过渡到强形式的关键。对 \(f\in k[x_1,\ldots,x_n]\)，再加变量 \(T\)，方程 \(Tf-1=0\) 正是要求 \(f\) 的值可逆。这个方法称为\term[rabinowitsch-jiqiao]{Rabinowitsch 技巧}[Rabinowitsch trick]。\footnote{原论文为《Zum Hilbertschen Nullstellensatz》，发表于 1930 年的《Mathematische Annalen》，见\cite{Rabinowitsch1930Nullstellensatz}。}

其中从增广理想回到根理想的推理，不依赖零点或代数闭性，而是下面的纯代数事实。

\begin{lemma}\label{b3:lem:rabinowitsch-equivalence}
设 \(R\) 为交换含幺环，\(I\subseteq R\) 为理想，\(f\in R\)，\(T\) 为不定元。则
\[
f\in\sqrt I
\quad\Longleftrightarrow\quad
IR[T]+(1-Tf)=R[T].
\]
\end{lemma}
\begin{proof}
若 \(f^N\in I\)、\(N\ge1\)，有限几何级数给出
\[
1=(1-Tf)\sum_{j=0}^{N-1}(Tf)^j+T^Nf^N,
\]
所以右侧理想含一。

反向，设右侧理想含一。由多项式环的泛性质，局部化映射 \(R\to R_f\) 延拓为环同态
\[
R[T]\longrightarrow R_f,\qquad T\longmapsto f^{-1}.
\]
在这个同态下，\(1-Tf\) 变成零，故单位理想的等式给出 \(1\in IR_f\)。将\cref{b3:prop:ideal-localization-saturation}用于乘法集 \(\{1,f,f^2,\ldots\}\)，得到 \(I\) 中含有某个 \(f^N\)。若 \(N=0\)，则 \(1\in I\)，可改取 \(N=1\)。因此 \(f\in\sqrt I\)。这里通过局部化的成员判据清分母，所以即使 \(R\) 有零因子，推理仍成立。
\end{proof}

\begin{proof}[\cref{b3:thm:strong-nullstellensatz}的证明]
若 \(f^r\in J\)，则在 \(Z_k(J)\) 的每个点都有 \(f(a)^r=0\)，从而 \(f(a)=0\)。因此 \(\sqrt J\subseteq I\bigl(Z_k(J)\bigr)\)。

反向，取在 \(Z_k(J)\) 上恒为零的多项式 \(f\)，在 \(k[x_1,\ldots,x_n,T]\) 中考虑
\[
J' = J\,k[x_1,\ldots,x_n,T]+(1-Tf).
\]
若 \((a,t)\) 为其零点，则 \(a\in Z_k(J)\)，所以 \(f(a)=0\)，但又要求 \(1-t\cdot f(a)=0\)，矛盾。由\cref{b3:cor:proper-ideal-common-zero}，无公共零点使 \(J'\) 成为单位理想；再由\cref{b3:lem:rabinowitsch-equivalence}，得到 \(f\in\sqrt J\)。弱零点定理在这里把几何上的无解转成单位理想，随后根理想的结论完全由纯代数引理给出。

两种运算都反向包含。刚证明对根理想 \(J\) 有 \(I\bigl(Z_k(J)\bigr)=J\)。对代数集 \(X=Z_k(J)\)，一方面 \(J\subseteq I(X)\) 给出 \(Z_k\bigl(I(X)\bigr)\subseteq X\)；另一方面每个 \(a\in X\) 按定义使 \(I(X)\) 中全部多项式为零，得到反向包含。于是两种运算在所述对象上互逆。
\end{proof}

\subsection{几何与代数的反变关系}\label{b3:subsec:0603:03}

\begin{definition}\label{b3:def:affine-coordinate-ring}
设 \(k\) 为代数闭域，\(n\ge0\)，\(X\subseteq k^n\) 为多项式方程组的公共零点集。定义商环
\[
k[X]=k[x_1,\ldots,x_n]/I(X)
\]
并称它为 \(X\) 的\term[zuobiao-huan]{坐标环}[coordinate ring]。其中一个元素就是多项式在 \(X\) 上给出的函数；两个多项式给出同一函数，当且仅当它们之差属于 \(I(X)\)。
\end{definition}
\symidx{\(k[X]\)：仿射代数集 \(X\) 的坐标环}

因为 \(I(X)\) 是根理想，普通代数集的坐标环总是约化环。若从理想 \(J\subseteq R=k[x_1,\ldots,x_n]\) 出发，先取零点集再取坐标环，强零点定理给出
\[
k[Z_k(J)]=R/\sqrt J.
\]
这未必是 \(R/J\)：例如前面的 \(J=(x^2,y)\) 得到原点的坐标环 \(k\)，而 \(R/J\) 中的非零幂零元 \(\bar x\) 已经丢失。

设 \(X,Y\subseteq k^n\) 为代数集。强零点定理说明，代数集的包含与零点理想的包含反向对应：
\[
X\subseteq Y\ \Longleftrightarrow\ I(Y)\subseteq I(X).
\]
当 \(X\subseteq Y\) 时，把 \(Y\) 上的多项式函数限制到 \(X\)，就得到满同态 \(k[Y]\rightarrow k[X]\)，其核是 \(I(X)/I(Y)\)。因此
\[
k[X]\cong k[Y]\big/\bigl(I(X)/I(Y)\bigr).
\]
例如 \(Y=k^2\)、\(X=Z_k(y)\) 时，限制函数把 \(k[x,y]\) 送到 \(k[x]\)，正是令 \(y=0\) 的商同态。代数集的包含由此表现为函数环的商，而不是子环。

\ProseParagraphBreak
推广到一般映射时，需要判断函数与映射复合后是否仍属于相应的函数环；在局部允许分母的情形，还要考察这些表达式在各点是否有定义。\secref{b3:sec:0701}与\secref{b3:sec:0702}以此研究代数集的拓扑、正则函数与态射。\secref{b3:sec:0604}的尖点、节点和\secref{b3:sec:0606}的导子计算，则进一步比较零点集与环本身保留的信息。
